\documentclass{article}
\usepackage[utf8]{inputenc}
\usepackage{amsmath}
\usepackage{amssymb}
\usepackage{graphicx}
\usepackage{tabularx}
\usepackage{booktabs}
\usepackage{caption}
\usepackage[english]{babel}
\usepackage{setspace}
\usepackage{geometry}
\usepackage{float}
\usepackage{url}
\usepackage[justification=centering]{caption}

% APA style formatting
\geometry{a4paper, margin=1in}
\setstretch{1.5}
\captionsetup[table]{labelfont=bf, justification=raggedright, singlelinecheck=false}
\captionsetup[figure]{font=small, labelfont=bf, justification=raggedright, singlelinecheck=false}

\title{ Machine Learning for Customer Churn Prediction}
\author{Nelson David Posso Suarez \\ Javier Alejandro Penagos Hernandez \\ Nicolas Castro Rivera }
\date{}

\begin{document}

\maketitle

\section{Summary}

This work formulates the customer churn retention problem as a Markov Decision Process (MDP).
The objective is to train a Reinforcement Learning (RL) agent to select the best retention strategy
for a customer with a probability of churn. Although the available dataset does not contain labeled
data describing the outcome of different retention strategies, it provides sufficient customer and
service information to establish an environment where the agent can learn which actions can reduce
the probability of churn.



\section{MDP Formulation}
The MDP is composed of a state space containing customer features, an action set representing
different retention strategies, stochastic transitions based on the customer's acceptance of an
offer, and a reward function based on the reduction of the predicted churn probability and the
cost of the selected action. RL agent will learn using DQN or PPO to determine the action with 
the highest expected reward. 

\subsection{Problem Statement}

The objective of the MDP is to determine the best retention strategy for a customer according to
their characteristics and current probability of churn. The RL agent receives the customer's state
and selects an action from a set of possible retention strategies. The selected action can be
accepted or rejected by the customer, introducing stochasticity into the environment.

The customer's churn probability is obtained using a supervised learning model. This probability
is used by the RL environment to evaluate whether the selected retention strategy improves the
customer's situation. The objective of the agent is therefore to reduce the probability of churn
while considering the cost associated with some retention strategies.

\subsection{State Space}

The state consist in the customer features, the data processed generated in the preprocessing pipeline. 
Any row can be a state valid for an episode in the environment and plus the probability of churn. 
\\There are 3 sections: member data of the user like age, city, etc. Transaction data calculated with the historical transactions.
User logs data calculated with logs in the last 3 months and probability churn. 
    
For this reason, each state is represented as a vector of 60 elements:

\begin{center}
\texttt{State = [Members 8 + transactions 29 +  user logs 22 + churn probability 1 = 60]}
\end{center}


\subsection{Action Set}

The action set are three discrete retention strategies.
Each action is designed to target a different group of features in the preprocessed state vector,
so that no single strategy dominates all customer profiles and the agent is encouraged to learn
context-dependent selection. Actions are validated against the customer's current subscription
status before being offered, and each is subject to a stochastic acceptance by the customer.

\subsubsection{Renewal Offer}

Customers whose subscription has lapsed and have cancelled their
most recent transaction or disabled auto-renewal has a big chance of be churn.
\\This strategy offer to renew and is only valid when the customer's last
transaction satisfies \texttt{is\_cancel = 1} or \texttt{is\_auto\_renew = 0}.
\\The acceptance probability is 40\%. This reflects the difficulty of winning back a customer who
has already taken a deliberate step away from the service. The action has no economic cost for the company.

When accepted, the environment injects a new monthly transaction record into the customer's raw
data, restoring \texttt{is\_auto\_renew = 1} and \texttt{is\_cancel = 0} while preserving the
existing plan duration and price. The features most directly affected are
\texttt{last\_is\_cancel}, \texttt{last\_is\_auto\_renew}, \texttt{days\_to\_expire}, and
\texttt{days\_since\_last\_transaction}.



\subsubsection{Annual Plan and Re-engagement}

This strategy targets customers not yet on an annual plan, specifically those whose most recent
transaction has \texttt{payment\_plan\_days} $< 365$. The offer upgrades the customer to a
365-day plan at a list price of 1,490 NTD with a built-in 20\% discount, so the effective
amount paid is 1,192 NTD.

Beyond the new transaction, the action also injects four synthetic listening-session records
dated in the final days of the feature window. These simulate the re-engagement behaviour that
typically follows a successful annual plan offer and propagate into the listening trend features
(\texttt{trend\_songs}, \texttt{trend\_secs}, \texttt{trend\_ratio100}, \texttt{trend\_ratio25})
and active-day counts.

The acceptance probability is 70\%, is high for the discount but not enough for some customers that 
not want has the commitment with the company. The action carries a cost of 10\% because the company
absorbs the 20\% price reduction on the annual plan.

\subsubsection{Discount Offer}

This strategy offers a 5\% discount on the customer's next monthly renewal. It is only valid when
the customer is not already paying below list price
(\texttt{actual\_amount\_paid} $\geq$ \texttt{plan\_list\_price}). When accepted, the new
transaction records \texttt{actual\_amount\_paid} at 95\% of the current
\texttt{plan\_list\_price}. The features most directly affected are \texttt{last\_discount},
\texttt{average\_amount\_paid}, and \texttt{difference}.

The acceptance probability is \textbf{dynamic}: 85\% for active subscribers and 20\% for
lapsed or cancelled customers. A price reduction is compelling for a subscriber still engaged
with the service, but has little effect on one who has already decided to leave. The action
carries a cost of 15\%.

\subsection{Transition Dynamics}

The selected action directly affects the customer's state. However, before modifying the state,
the environment determines whether the customer accepts or rejects the offered strategy. This
makes the transition stochastic because the same action can produce different outcomes.

The acceptance probabilities are defined according to the type of retention strategy:

\begin{itemize}
\item Renewal Offer: 40\% acceptance and 60\% rejection.
\item Annual Plan and Re-engagement: 70\% acceptance and 30\% rejection.
\item Discount Offer: 85\% or 20\% acceptance and 15\% or 80\% rejection respectively.
\end{itemize}

The transition function can be represented as:

\begin{equation}
s_{t+1} = f(s_t, a_t)
\end{equation}

The acceptance or rejection of an action is stochastic. However, once an action is accepted,
the modification applied to the customer's state is deterministic. Therefore, the same state
and the same accepted action produce the same resulting state.

\subsection{Reward Function}

The reward is the reduction in predicted churn probability discounted by the action cost, 
the agent search maximize the reward, if the reduction difference of probabilities is high the reward will be increase.
But some strategies can generate cost to company and can penalize the use of this for retain de the customer.  

\begin{equation}
r_t = \bigl(p_{\text{before}} - p_{\text{after}}\bigr)\,(1 - c_{a_t})
\end{equation}

where $p_{\text{before}}$ and $p_{\text{after}}$ are the supervised model's churn probability
estimates before and after the action is applied, and $c_{a_t} \in \{0,\, 0.10,\, 0.15\}$ is
the cost of the selected action. When an action is rejected the state is unchanged and
$r_t = 0$.

\subsection{Episode }

An episode starts with a customer from the KKBox dataset. The customer's initial
state is provided to the supervised learning model, which predicts the probability of churn.
This information is then provided to the RL agent.

The agent select the retention strategy with the highest expected reward.
The selected action is then evaluated by the environment, which determines whether the customer
accepts or rejects the offer according to the corresponding acceptance probability.

If the customer accepts the offer, the action modifies the customer's state. The supervised model
then predicts the new probability of churn, which is compared with the probability before the action
to calculate the reward. The episode terminates after this reward is obtained.

If the customer rejects the offer, the customer's state remains unchanged and the reward is zero.
The episode then terminates.


\section{Environment}

\section{Agent algorithm}

\section{Hyperparameter}

\section{Sample trajectory visualizations}

\section{Discussion}

\section{Conclusion}


\begin{thebibliography}{9}

\bibitem{ahmad2019}
Ahmad, A. K., Jafar, A., \& Aljoumaa, K. (2019). Customer churn prediction in telecom using machine learning in big data platform. \textit{Journal of Big Data}, 6, 28.

\bibitem{verbeke2012}
Verbeke, W., Dejaeger, K., Martens, D., Hur, J., \& Baesens, B. (2012). New insights into churn prediction in the telecommunication sector: A profit driven data mining approach. \textit{European Journal of Operational Research}, 218(1), 211--229.

\bibitem{breiman2001randomforests}
Breiman, L. (2001). Random forests. \textit{Machine Learning}, 45(1), 5--32.

\bibitem{chen2016xgboost}
Chen, T., \& Guestrin, C. (2016). XGBoost: A scalable tree boosting system. In \textit{Proceedings of the 22nd ACM SIGKDD International Conference on Knowledge Discovery and Data Mining} (pp. 785--794).

\bibitem{akiba2019optuna}
Akiba, T., Sano, S., Yanase, T., Ohta, T., \& Koyama, M. (2019). Optuna: A next-generation hyperparameter optimization framework. In \textit{Proceedings of the 25th ACM SIGKDD International Conference on Knowledge Discovery \& Data Mining} (pp. 2623--2631).

\bibitem{paszke2019pytorch}
Paszke, A., Gross, S., Massa, F., Lerer, A., Bradbury, J., Chanan, G., Killeen, T., Lin, Z., Gimelshein, N., Antiga, L., Desmaison, A., Kopf, A., Yang, E., DeVito, Z., Raison, M., Tejani, A., Chilamkurthy, S., Steiner, B., Fang, L., Bai, J., \& Chintala, S. (2019). PyTorch: An imperative style, high-performance deep learning library. \textit{Advances in Neural Information Processing Systems}, 32.

\bibitem{chawla2002smote}
Chawla, N. V., Bowyer, K. W., Hall, L. O., \& Kegelmeyer, W. P. (2002). SMOTE: Synthetic minority over-sampling technique. \textit{Journal of Artificial Intelligence Research}, 16, 321--357.

\bibitem{davis2006precision}
Davis, J., \& Goadrich, M. (2006). The relationship between precision-recall and ROC curves. In \textit{Proceedings of the 23rd International Conference on Machine Learning} (pp. 233--240).

\end{thebibliography}

\end{document}
